\documentclass[10pt,a4paper]{amsart}
\usepackage[margin=19mm]{geometry}
\usepackage{amsmath,amssymb,mathtools}
\usepackage[scr=rsfs,cal=euler]{mathalfa}
\usepackage{xfrac}
\usepackage[unicode,hidelinks,bookmarks=false]{hyperref}
\newcommand{\ie}{i.e.\ }
\newcommand{\cf}{cf.\ }
\newcommand{\resp}{respectively\ }
\newcommand{\Subsection}{\subsection}
\providecommand{\nobreakdash}{\nobreak\hbox{-}}
\setlength{\emergencystretch}{2em}
\multlinegap0pt
\usepackage{amssymb}
\usepackage{float}
\usepackage{tikz}
\newtheorem{thm}{Theorem}
\title{\textbf{Representability of Flag Matroids}}
\author{Daniel Irving Bernstein, Nathaniel Vaduthala}
\date{Source: arXiv:2504.00267v2 (CC BY 4.0)}
\begin{document}
\maketitle

\noindent\emph{Source and licence.} \textbf{Representability of Flag Matroids}. Version arXiv:2504.00267v2 (CC BY 4.0), distributed by arXiv under the Creative Commons Attribution 4.0 International licence, http://creativecommons.org/licenses/by/4.0/, as recorded in the arXiv licence metadata for this exact version and shown on the abstract page https://arxiv.org/abs/2504.00267v2. Full licence notice: \texttt{licenses/arxiv-2504.00267-CC-BY-4.0.txt}.\par\smallskip

\noindent\textit{Editorial scope.} Counted unit: the article's thm representability (source0.tex lines 439--441). The article's complete original proof of it is delivered with it (source0.tex lines 442--464, one \texttt{proof} environment of 23 lines). No mathematics is written, repaired or re-derived: the statement and the proof are the article's own lines, in the article's own notation, and no retained line is substituted or re-flowed.  No further source-local context is needed: every cross-reference of every retained line resolves inside the two delivered spans.  Nothing was omitted from any retained span, so the package carries no omission record.  No retained line contains a citation, so the wrapper carries no bibliography and no bibliography entry.  No retained line contains a figure environment, an image-inclusion command or drawing code, so no image is delivered and the package declares no asset dependency.  Editorial formatting only, in addition: an AMS \texttt{amsart} preamble that restates the article's own declarations and macros used by the retained lines, namely the environments \texttt{thm} on separate counters, together with the standard packages those lines need.  The retained environments print in this document as Theorem 1 in order of appearance, whereas the article prints the same items with the numbering of its own full document.  The complete native source of the article as distributed in the arXiv e-print for this exact version is bundled unchanged as \texttt{sources/175735/source0.tex}.
\begin{thm}\label{thm: uniform flag matroids representability}
    The flag matroid $\mathcal{I}(U_{r,n})$ is $\mathbb{K}$-representable if and only if $r\leq 1$ or $|\mathbb{K}| \geq n$.
\end{thm}

\begin{proof}
    The constant matrix of size $1\times n$, with each entry equal to $0$ (respectively $1$) is a representation of $U_{0,n}$ (respectively $U_{1,n}$) over any field.
    Suppose now $|\mathbb{K}| \geq n$.
    Let $e_1, e_2, \dots, e_n$ be distinct elements of $\mathbb{K}$ and define 
    \begin{equation*}
        A := \begin{pmatrix}
            1 & 1 & \cdots & 1 \\
            e_1 & e_2 & \cdots & e_n \\
            e_1^2 & e_2^2 & \cdots & e_n^2 \\
            \vdots & \vdots & \ddots & \vdots \\
            e_1^{r-1} & e_2^{r-1} & \dots & e_n^{r-1}
        \end{pmatrix}
    \end{equation*}
    If $F \subseteq \{1,\dots,n\}$ with $|F| \le r$,
    then $A_{1, 2, \dots, |F|}^F$ is a Vandermonde matrix and therefore nonsingular.
    Since $A$ has size $r\times n$, this implies that $\mathfrak{F}(A) = \mathcal{I}(U_{r,n})$.

    Now assume $\mathcal{I}(U_{r,n})$ is $\mathbb{K}$-representable and suppose $r > 1$.
    The first row of $A$ cannot contain any zero entries
    so by appropriately scaling each column we may assume that the first row of $A$ is all ones.
    Since all subsets of size 2 are feasible in $\mathcal{I}(U_{r,n})$,
    all entries of the second row of $A$ must be distinct and so $|\mathbb{K}| \ge n$.
\end{proof}

\end{document}
